% Fichier généré automatiquement par tools/generate_corrections.py.
% Source : CIN/CIN-02-VitesseAcceleration/02_R_02/02_R_02.tex
% Statut : complet (4/4 question(s) renseignée(s), 0 reprise(s)).
\begin{corrigebox}[Corrigé extrait de la banque]
\CorrigeQuestion{1}
$\vectv{B}{1}{0}$ $ = \deriv{\vect{AB}}{\rep{0}}$ $=\deriv{R\vect{i_1}}{\rep{0}}$.
Or $\deriv{\vi{1}}{\rep{0}} = \deriv{\vi{1}}{\rep{1}} + \vecto{1}{0}\wedge \vi{1} $ $=\vect{0}+\dot{\theta}\vk{0}\wedge \vi{1}$ $=\dot{\theta}\vj{1}$.

D'où $\vectv{B}{1}{0}=R\dot{\theta}\vj{1}$.
\CorrigeQuestion{2}
$\babarv{B}{A}{1}{0}$ $=\vect{0}-R\vi{1}\wedge \dot{\theta}\vk{0} =R\dot{\theta}\vj{1} $.
\CorrigeQuestion{3}
On a directement $\torseurcin{V}{1}{0} = \torseurl{\dot{\theta}\vk{0}}{R\dot{\theta}\vj{1}}{B}$.
\CorrigeQuestion{4}
$\vectg{B}{1}{0} = \deriv{\vectv{B}{1}{0}}{\rep{0}}=R\ddot{\theta}\vj{1}-R\dot{\theta}^2\vi{1}$.
(En effet,  $\deriv{\vj{1}}{\rep{0}} = \deriv{\vj{1}}{\rep{1}} + \vecto{1}{0}\wedge \vj{1} $ $=\vect{0}+\dot{\theta}\vk{0}\wedge \vj{1}$ $=-\dot{\theta}\vi{1}$.)
\end{corrigebox}
